\chapter{General Relativity Background}
\label{chap:general-relativity}

This chapter summarizes the parts of general relativity that are needed in the
rest of the thesis. We work in geometrized units \(G=c=1\) and use the metric
signature \((-+++)\). The systems studied below are test-particle systems: the
particle moves in a prescribed space-time geometry and its own gravity is
neglected. The relevant geometry is therefore encoded in the metric, its
symmetries, and the geodesic equation.

\section{Geodesic Motion}
\label{sec:geodesic-motion}

A space-time is a four-dimensional Lorentzian manifold with metric tensor
\(g_{\mu\nu}\). In local coordinates \(x^\mu\), the line element is
\begin{equation}
    \mathrm{d}s^2
    =
    g_{\mu\nu}\,\mathrm{d}x^\mu\mathrm{d}x^\nu .
    \label{eq:line-element}
\end{equation}
For a massive free test particle, the world-line can be parametrized by its
proper time \(\tau\). Its four-velocity is
\begin{equation}
    u^\mu
    =
    \frac{\mathrm{d}x^\mu}{\mathrm{d}\tau},
    \label{eq:four-velocity}
\end{equation}
and the normalization of a timelike world-line is
\begin{equation}
    g_{\mu\nu}u^\mu u^\nu=-1.
    \label{eq:timelike-normalization-chap2}
\end{equation}
The free motion is geodesic motion. In coordinates, the geodesic equation is
\begin{equation}
    \frac{\mathrm{d}^2x^\mu}{\mathrm{d}\tau^2}
    +
    \Gamma^\mu_{\alpha\beta}
    \frac{\mathrm{d}x^\alpha}{\mathrm{d}\tau}
    \frac{\mathrm{d}x^\beta}{\mathrm{d}\tau}
    =
    0,
    \label{eq:geodesic-equation}
\end{equation}
where the Christoffel symbols are determined by the metric,
\begin{equation}
    \Gamma^\mu_{\alpha\beta}
    =
    \frac{1}{2}g^{\mu\nu}
    \left(
        g_{\nu\alpha,\beta}
        +
        g_{\nu\beta,\alpha}
        -
        g_{\alpha\beta,\nu}
    \right).
    \label{eq:christoffel-symbols}
\end{equation}
Thus, once the metric is known, the motion of a free test particle is fixed.
This is the basic geometric formulation of free fall in general relativity
\cite{BicakSemerak2026,Kraus2025}.

\section{Symmetries and Conserved Quantities}
\label{sec:symmetries-conserved-quantities}

The space-times used in this thesis are stationary and axially symmetric. They
therefore admit two commuting Killing vector fields: one associated with time
translations and one with rotations around the symmetry axis. In adapted
coordinates we denote them by
\begin{equation}
    \xi^\mu_{(t)}=\delta^\mu_t,
    \qquad
    \xi^\mu_{(\phi)}=\delta^\mu_\phi .
\end{equation}
If \(\xi^\mu\) is a Killing vector field, then
\(g_{\mu\nu}u^\mu\xi^\nu\) is conserved along a geodesic. For the two
symmetries above this gives the specific energy and azimuthal angular momentum
\begin{equation}
    \epsilon
    :=
    -g_{\mu\nu}u^\mu\xi^\nu_{(t)},
    \qquad
    \ell
    :=
    g_{\mu\nu}u^\mu\xi^\nu_{(\phi)}.
    \label{eq:killings}
\end{equation}
For a diagonal stationary-axisymmetric metric this means
\begin{equation}
    u^t=-\frac{\epsilon}{g_{tt}},
    \qquad
    u^\phi=\frac{\ell}{g_{\phi\phi}}.
    \label{eq:ut-uphi-conserved}
\end{equation}
The two cyclic coordinates \(t\) and \(\phi\) can therefore be eliminated from
the dynamical analysis. The nontrivial motion is reduced to the meridional
plane, described in this thesis by Weyl-type coordinates \((\rho,z)\).

Substituting Eq.~\eqref{eq:ut-uphi-conserved} into the timelike
normalization~\eqref{eq:timelike-normalization-chap2} gives
\begin{equation}
    -1
    =
    \epsilon^2 g^{tt}
    +
    \ell^2 g^{\phi\phi}
    +
    g_{\rho\rho}(u^\rho)^2
    +
    g_{zz}(u^z)^2 .
    \label{eq:reduced-normalization-chap2}
\end{equation}
This relation is the starting point for the accessible-region construction in
Chapter~\ref{chap:implementation}.

\section{Weyl Space-times}
\label{sec:weyl-spacetimes}

Static axisymmetric vacuum space-times can be written in Weyl coordinates
\((t,\phi,\rho,z)\). In the convention used here, their metric is
\begin{equation}
    \mathrm{d}s^2
    =
    -e^{2\nu}\,\mathrm{d}t^2
    +
    \rho^2 e^{-2\nu}\,\mathrm{d}\phi^2
    +
    e^{2(\lambda-\nu)}
    \left(
        \mathrm{d}\rho^2+\mathrm{d}z^2
    \right),
    \label{eq:weyl-metric}
\end{equation}
where \(\nu=\nu(\rho,z)\) and \(\lambda=\lambda(\rho,z)\). Equivalently,
with \(N=e^\nu\),
\begin{equation}
    \mathrm{d}s^2
    =
    -N^2\,\mathrm{d}t^2
    +
    N^{-2}
    \left[
        \rho^2\,\mathrm{d}\phi^2
        +
        e^{2\lambda}
        \left(
            \mathrm{d}\rho^2+\mathrm{d}z^2
        \right)
    \right].
    \label{eq:weyl-metric-lapse}
\end{equation}
The first metric function satisfies the flat-space axisymmetric Laplace
equation in vacuum,
\begin{equation}
    \nu_{,\rho\rho}
    +
    \frac{1}{\rho}\nu_{,\rho}
    +
    \nu_{,zz}
    =
    0.
    \label{eq:weyl-laplace}
\end{equation}
This is why \(\nu\) is often called the Weyl potential. In the weak-field
limit, \(g_{tt}=-e^{2\nu}\simeq -(1+2\nu)\), so \(\nu\) has the role of the
Newtonian gravitational potential. The second metric function is fixed by
\begin{align}
    \lambda_{,\rho}
    &=
    \rho
    \left[
        \left(\nu_{,\rho}\right)^2
        -
        \left(\nu_{,z}\right)^2
    \right],
    \label{eq:weyl-lambda-rho}
    \\
    \lambda_{,z}
    &=
    2\rho\nu_{,\rho}\nu_{,z}.
    \label{eq:weyl-lambda-z}
\end{align}
The equations for \(\lambda\) are nonlinear in \(\nu\). Therefore, Weyl
potentials can be added linearly, but the corresponding \(\lambda\) of a
superposed configuration contains interaction terms. This feature is important
for perturbed black holes: the first metric function behaves in a Newtonian-like
way, while the full geometry remains genuinely relativistic
\cite{Sychrovsky2020,Kraus2025,Semerak2016}.

\section{Schwarzschild Black Hole in Weyl Coordinates}
\label{sec:schwarzschild-weyl}

The Schwarzschild metric in standard coordinates \((t,r,\theta,\phi)\) is
\begin{equation}
    \mathrm{d}s^2
    =
    -\left(1-\frac{2M}{r}\right)\mathrm{d}t^2
    +
    \left(1-\frac{2M}{r}\right)^{-1}\mathrm{d}r^2
    +
    r^2
    \left(
        \mathrm{d}\theta^2
        +
        \sin^2\theta\,\mathrm{d}\phi^2
    \right).
    \label{eq:schwarzschild-metric}
\end{equation}
It describes the exterior of a nonrotating black hole of mass \(M\). The
transformation to Weyl coordinates is
\begin{equation}
    \rho
    =
    \sqrt{r(r-2M)}\sin\theta,
    \qquad
    z
    =
    (r-M)\cos\theta.
    \label{eq:schwarzschild-to-weyl}
\end{equation}
If
\begin{equation}
    d_{\pm}
    :=
    \sqrt{\rho^2+(z\pm M)^2},
    \label{eq:schwarzschild-distances}
\end{equation}
then
\begin{equation}
    r
    =
    M+\frac{d_{+}+d_{-}}{2},
    \qquad
    \cos\theta
    =
    \frac{d_{+}-d_{-}}{2M}.
    \label{eq:weyl-to-schwarzschild}
\end{equation}
The Schwarzschild horizon is represented by the finite rod
\begin{equation}
    \rho=0,
    \qquad
    -M\leq z\leq M.
\end{equation}
The Weyl functions of the Schwarzschild solution are
\begin{align}
    \nu_{\mathrm{Schw}}
    &=
    \frac{1}{2}
    \ln
    \left(
        \frac{d_{+}+d_{-}-2M}
             {d_{+}+d_{-}+2M}
    \right),
    \label{eq:schwarzschild-potential}
    \\
    \lambda_{\mathrm{Schw}}
    &=
    \frac{1}{2}
    \ln
    \left[
        \frac{(d_{+}+d_{-})^2-4M^2}
             {4d_{+}d_{-}}
    \right].
    \label{eq:schwarzschild-lambda}
\end{align}
These formulas make the Schwarzschild black hole directly compatible with the
Weyl superposition method used for the external ring
\cite{Sychrovsky2020,Kraus2025}.

\section{Bach--Weyl Ring and the Schwarzschild--Bach--Weyl System}
\label{sec:bach-weyl-ring}

The Bach--Weyl solution represents a static, axially symmetric,
infinitesimally thin circular ring. We denote by \(b\) the Schwarzschild
radial coordinate of the ring used for comparison with the literature. Its Weyl
cylindrical radius is
\begin{equation}
    \sigma_{\mathrm{BW}}
    =
    \sqrt{b(b-2M)},
    \label{eq:schw-bw-ring-radius}
\end{equation}
which follows from the Schwarzschild--Weyl transformation
\eqref{eq:schwarzschild-to-weyl} at \(\theta=\pi/2\). In the ring potential
below, \(\sigma\) denotes this Weyl radius. Let \(\mathcal{M}\) be the ring
mass and define
\begin{equation}
    l_{1,2}
    :=
    \sqrt{(\rho\mp \sigma)^2+z^2},
    \qquad
    k^2
    :=
    1-\frac{l_1^2}{l_2^2}
    =
    \frac{4\sigma\rho}{l_2^2}.
    \label{eq:ring-quantities}
\end{equation}
The complete elliptic integral of the first kind is
\begin{equation}
    K(k)
    :=
    \int_0^{\pi/2}
    \frac{\mathrm{d}\alpha}
         {\sqrt{1-k^2\sin^2\alpha}}.
    \label{eq:complete-elliptic-k}
\end{equation}
The Bach--Weyl potential is
\begin{equation}
    \nu_{\mathrm{BW}}
    =
    -\frac{2\mathcal{M}K(k)}{\pi l_2}.
    \label{eq:bw-potential}
\end{equation}
It is the relativistic counterpart of the Newtonian potential of a homogeneous
thin ring. The function \(\lambda_{\mathrm{BW}}\), however, is not Newtonian.
It is obtained from Eqs.~\eqref{eq:weyl-lambda-rho} and
\eqref{eq:weyl-lambda-z}; near the ring it is direction-dependent and singular.
This makes the Bach--Weyl ring geometrically more severe than its Newtonian
analogue \cite{Semerak2016,Sychrovsky2020,Kraus2025}.

Since Eq.~\eqref{eq:weyl-laplace} is linear, the Schwarzschild and
Bach--Weyl potentials can be superposed:
\begin{equation}
    \nu
    =
    \nu_{\mathrm{Schw}}
    +
    \nu_{\mathrm{BW}}.
    \label{eq:schwarzschild-bw-potential}
\end{equation}
The lapse is consequently
\begin{equation}
    N
    =
    e^\nu
    =
    N_{\mathrm{Schw}}N_{\mathrm{BW}}.
    \label{eq:schwarzschild-bw-lapse}
\end{equation}
The second Weyl function is not a simple sum. It has the structure
\begin{equation}
    \lambda
    =
    \lambda_{\mathrm{Schw}}
    +
    \lambda_{\mathrm{BW}}
    +
    \lambda_{\mathrm{int}},
    \label{eq:schwarzschild-bw-lambda}
\end{equation}
where \(\lambda_{\mathrm{int}}\) is the interaction contribution generated by
the quadratic equations for \(\lambda\). The combined metric
Eqs.~\eqref{eq:weyl-metric} and
\eqref{eq:schwarzschild-bw-potential} is the Schwarzschild--Bach--Weyl
(SCHW+BW) space-time used in the first part of the numerical study. In
the following chapters, the abbreviation SCHW+BW always refers to this
Schwarzschild black hole plus Bach--Weyl ring system; SCHW denotes the
Schwarzschild black hole and BW denotes the Bach--Weyl ring.

\section{Majumdar--Papapetrou Space-times}
\label{sec:majumdar-papapetrou}

Majumdar--Papapetrou space-times are static electrovacuum solutions in which
the gravitational attraction of extremally charged sources is balanced by their
electrostatic repulsion. In cylindrical coordinates they can be written as
\begin{equation}
    \mathrm{d}s^2
    =
    -N^2\,\mathrm{d}t^2
    +
    N^{-2}
    \left(
        \rho^2\,\mathrm{d}\phi^2
        +
        \mathrm{d}\rho^2
        +
        \mathrm{d}z^2
    \right).
    \label{eq:mp-metric}
\end{equation}
This is a special case of the Weyl form with \(\lambda=0\). The practical
variable is \(N^{-1}\), which is harmonic in flat three-dimensional space away
from the sources. For several sources one has a linear superposition in
\(N^{-1}\), rather than in \(\nu=\ln N\) \cite{Kraus2025,Polcar2019}.

\subsection{Extreme Reissner--Nordstr\"om Black Hole}
\label{subsec:extreme-rn}

The Reissner--Nordstr\"om metric in spherical coordinates is
\begin{equation}
    \mathrm{d}s^2
    =
    -f(r)\,\mathrm{d}t^2
    +
    f(r)^{-1}\,\mathrm{d}r^2
    +
    r^2
    \left(
        \mathrm{d}\theta^2
        +
        \sin^2\theta\,\mathrm{d}\phi^2
    \right),
    \label{eq:rn-metric}
\end{equation}
where
\begin{equation}
    f(r)
    =
    1-\frac{2M}{r}+\frac{Q^2}{r^2}.
    \label{eq:rn-function}
\end{equation}
For the extreme case, \(|Q|=M\), one obtains
\begin{equation}
    f(r)
    =
    \left(1-\frac{M}{r}\right)^2.
\end{equation}
The corresponding Weyl-type coordinates satisfy
\begin{equation}
    \rho=(r-M)\sin\theta,
    \qquad
    z=(r-M)\cos\theta.
    \label{eq:ern-weyl-coordinates}
\end{equation}
Thus \(R=\sqrt{\rho^2+z^2}=r-M\), and the extreme
Reissner--Nordstr\"om lapse can be written as
\begin{equation}
    \frac{1}{N_{\mathrm{eRN}}}
    =
    1+\frac{M}{\sqrt{\rho^2+z^2}}.
    \label{eq:ern-lapse}
\end{equation}
This shows explicitly that the extreme Reissner--Nordstr\"om black hole belongs
to the Majumdar--Papapetrou class \cite{BicakSemerak2026,Polcar2019}.

\subsection{Majumdar--Papapetrou Ring and the Reissner--Nordstr\"om--Majumdar--Papapetrou System}
\label{subsec:mp-ring}

The Majumdar--Papapetrou ring is the extremally charged analogue of a thin
homogeneous circular ring. We again use \(b\) for the radial coordinate used in
the corresponding Reissner--Nordstr\"om coordinates. In the extreme
Reissner--Nordstr\"om case, Eq.~\eqref{eq:ern-weyl-coordinates} gives
\(R=\sqrt{\rho^2+z^2}=r-M\). Comparing the standard extreme
Reissner--Nordstr\"om factor \(g_{tt}=-(1-M/r)^2\) with the
Majumdar--Papapetrou form \(g_{tt}=-N^2\), \(N^{-1}=1+M/R\), gives
\(r=R+M\). Thus a ring at \(r=b\) is represented in the
Majumdar--Papapetrou cylindrical coordinates by
\begin{equation}
    \sigma_{\mathrm{MP}}
    =
    b-M.
    \label{eq:rn-mp-ring-radius}
\end{equation}
For a ring of cylindrical radius \(\sigma\) and mass \(\mathcal{M}\), the
inverse lapse is
\begin{equation}
    \frac{1}{N_{\mathrm{MP}}}
    =
    1
    +
    \frac{2\mathcal{M}K(k)}{\pi l_2},
    \label{eq:mp-ring-lapse}
\end{equation}
with \(l_1,l_2\) and \(k\) defined by Eq.~\eqref{eq:ring-quantities}, using
\(\sigma=\sigma_{\mathrm{MP}}\). In
contrast to the Bach--Weyl ring, the Majumdar--Papapetrou ring has
\(\lambda=0\). Its near-ring geometry is therefore less directional, which is
one reason it is useful as a comparison system \cite{Semerak2016,Polcar2019}.

Because both the extreme Reissner--Nordstr\"om black hole and the
Majumdar--Papapetrou ring belong to the same class of solutions, the combined
Reissner--Nordstr\"om--Majumdar--Papapetrou (RN+MP) system is obtained by
adding their contributions to \(N^{-1}\):
\begin{equation}
    \frac{1}{N}
    =
    1
    +
    \frac{M}{\sqrt{\rho^2+z^2}}
    +
    \frac{2\mathcal{M}K(k)}{\pi l_2}.
    \label{eq:ern-mp-superposition}
\end{equation}
Together with Eq.~\eqref{eq:mp-metric}, this gives the RN+MP metric used in
Chapter~\ref{chap:results}. In the following chapters, RN+MP denotes the
extreme Reissner--Nordstr\"om black hole plus Majumdar--Papapetrou ring
system; RN denotes Reissner--Nordstr\"om and MP denotes
Majumdar--Papapetrou.

\section{Effective Potential and Accessible Motion}
\label{sec:effective-potential}

For any diagonal stationary-axisymmetric metric of the form used above,
Eq.~\eqref{eq:reduced-normalization-chap2} can be rearranged as
\begin{equation}
    g_{\rho\rho}(u^\rho)^2+g_{zz}(u^z)^2
    =
    \epsilon^2\left(-g^{tt}\right)
    -
    \left(1+\ell^2g^{\phi\phi}\right).
    \label{eq:meridional-kinetic-term}
\end{equation}
Equivalently, since \(g_{tt}<0\),
\begin{equation}
    \epsilon^2
    \geq
    V_{\mathrm{eff}},
    \qquad
    V_{\mathrm{eff}}
    :=
    -g_{tt}
    \left(
        1+\ell^2g^{\phi\phi}
    \right).
    \label{eq:effective-potential-general}
\end{equation}
The allowed part of the meridional plane is the region where the right-hand
side of Eq.~\eqref{eq:meridional-kinetic-term} is non-negative. This condition
is the relativistic analogue of an energy bound in a reduced mechanical system.
It is also the geometric reason why basin maps can be drawn in a two-dimensional
initial-condition plane.

For the Weyl metric~\eqref{eq:weyl-metric},
\begin{equation}
    V_{\mathrm{eff}}^{\mathrm{Weyl}}
    =
    e^{2\nu}
    \left(
        1+\frac{\ell^2e^{2\nu}}{\rho^2}
    \right).
    \label{eq:effective-potential-weyl}
\end{equation}
For the Majumdar--Papapetrou metric~\eqref{eq:mp-metric},
\begin{equation}
    V_{\mathrm{eff}}^{\mathrm{MP}}
    =
    N^2
    +
    \frac{\ell^2N^4}{\rho^2}.
    \label{eq:effective-potential-mp}
\end{equation}
These two expressions are the theoretical basis for the accessible-region
plots and for the choice of initial conditions in the basin-map construction.
They also explain why the two systems behave differently near their central
objects: in the Schwarzschild--Bach--Weyl case the Weyl representation
contains the Schwarzschild horizon rod, whereas in the
Reissner--Nordstr\"om--Majumdar--Papapetrou case the centrifugal term and the
Majumdar--Papapetrou lapse can close or open the accessible region depending on
how the meridional plane is sampled.
